\section{Discussion: the open problems}\label{sec:discussion}
%%% ====================================================================

Elvang et al.~\cite{Elvang2026} show that in the $\mathcal{N}=4$ SYM EFT:
the Wilson coefficients $a_{2k,q}$ are fixed by SUSY and factorisation,
but $a_{2k-1,0} = \alpha^{\prime\,2k+1}\zeta(2k+1)$ remain free---they
correspond to odd Riemann zeta values which, from the EFT alone, are
external data.  The unique UV completion is the Veneziano amplitude.

The polylogarithmic circuit of \cref{ssec:polylog-circuit} identifies
the source of this apparent freedom: the EFT operates in the
\emph{additive} world of perturbative series ($\sum n^{-s}$), where the
$\zeta(2k+1)$ appear as irreducible constants.  The PCF arithmetic
operates in the \emph{multiplicative} fragment
($\mathbb{F}_1$-descent, Euler product $\prod f_p$), where the
same values are structurally determined by
$\golden = 2\cos(\pi/5)$.  The EFT has no access to this
multiplicative classification: it sees the values but not the
pentagonal mechanism that produces them.

Concretely:
\begin{itemize}
\item $\zeta(2k) = \pi^{2k}\cdot\mathbb{Q}$ (Euler formula): fixed by $\pi$
  alone, visible in the EFT via $a_{0,0}=\zeta(2)=\pi^2/6$.
\item $\zeta(2k+1) = \zeta_{\mathbb{Q}(\sqrt{5})}(2k+1)/L(2k+1,\chi_5)$: both factors are
  $\chi_5(a)$-projections of Hurwitz at $\{a/5\}$; the weights
  $\chi_5(a)=\log|2\cos(a\pi/5)|/\log\golden$ are read from the pentagon.
  These are not free---they are observations of the torus.
\item The Veneziano amplitude is the unique UV completion of the EFT,
  \emph{and} its odd-zeta coefficients are determined by the pentagon
  through the Frobenius structure $\psi_p(\golden)=\golden^p$.
  Elvang's EFT and PCF are complementary: the EFT identifies the
  \emph{which} (Veneziano), PCF explains the \emph{why} (pentagonal arithmetic).
\end{itemize}

The Euler product is the common object linking the three readings
of $T^2_{\mathrm{PCF}}$: physically, its period $T = 2\pi R$
governs the worldsheet partition function; arithmetically,
its factors are organised by the Frobenius structure
$\psi_p(\golden) = \golden^p$; categorically, they enter $\Ztwenty$
and determine the spectral structure through the Primitive Structure~\cite{V11}.
The pentagon generates all three.
The complementarity between EFT and PCF extends beyond the numerical values
of $\zeta(2k+1)$: the \emph{structures} that produce them---the Regge tower
on the string side, the Euler product on the arithmetic side---are the same
object viewed from two angles.

\smallskip\noindent\textbf{The Regge tower is the Euler product.}
The Veneziano amplitude~\cite{Veneziano}
\[
  A_4 = \frac{\Gamma(-\alpha' s)\,\Gamma(-\alpha' u)}
             {\Gamma(1-\alpha'(s+u))}
\]
has poles at $\alpha' s = n$ ($n = 0,1,2,\ldots$), forming the Regge tower
of string states at mass level~$n$.  Expanding in $\alpha'$, the residues
assemble the Dirichlet series $\sum_{n=1}^{\infty} n^{-s} = \zeta(s)$.
The fundamental theorem of arithmetic factorises each level
$n = p_1^{a_1}\cdots p_r^{a_r}$, reorganising the sum into the Euler
product $\prod_p(1-p^{-s})^{-1}$.

This reorganisation has a physical interpretation: the Euler product is the
\emph{second quantisation} of the prime spectrum.  Each prime $p$ is a
single-particle state $|p\rangle$; each mass level $n$ is the multi-particle
state $|p_1^{a_1},\ldots,p_r^{a_r}\rangle$; and the partition function of
this gas of primes is
\[
  Z(s) = \prod_{p}(1-p^{-s})^{-1} = \zeta(s).
\]

\smallskip\noindent\textbf{The $\chi_5$ grading classifies the string modes.}
PCF adds the structure that the EFT cannot see.  The character $\chi_5$
classifies each prime state:
\begin{itemize}
\item \textbf{Split} ($\chi_5(p)=+1$): two single-particle modes,
  contributing $f_p(s) = (1-p^{-s})^{-2}$.
  The prime $p$ sees both ideals above it in
  $\mathbb{Z}[\golden]$: $(p)=\mathfrak{P}\bar{\mathfrak{P}}$.
\item \textbf{Inert} ($\chi_5(p)=-1$): one single-particle mode at doubled
  energy, contributing $f_p(s) = (1-p^{-2s})^{-1}$.
  The prime $p$ sees one fused ideal of norm~$p^2$.
\end{itemize}
By the Chebotarev density theorem, the proportion of split primes
converges to $1/2$---the
same value $\mu_3 = 1/2$ that the Primitive Structure~\cite{V11}
identifies as the norm of the spectral operator.
The ideal-counting function $d_K(n) = \sum_{d\mid n}\chi_5(d)$
determines the Dirichlet coefficients of $\zeta_K(s)$: each mass
level~$n$ inherits its $\chi_5$-type from the factorisation of~$n$
into classified primes.

\smallskip\noindent\textbf{The polylogarithm is the Fourier analysis
of the Regge tower.}
The polylogarithm at a fifth root of unity,
\[
  \mathrm{Li}_s(\omega^a)
  = \sum_{n=1}^{\infty}\frac{e^{2\pi i n a/5}}{n^s},
\]
is the discrete Fourier transform of the Regge tower $\{1/n^s\}$ over
$\mathbb{Z}/5\mathbb{Z}$.
The $\chi_5$-projection extracts
$L(s,\chi_5) = (1/\sqrt{5})\sum_{a=1}^{4}\chi_5(a)\,
\mathrm{Li}_s(\omega^a)$
(\cref{ssec:polylog-circuit}).
Levels with $\chi_5(n)=+1$ contribute constructively;
levels with $\chi_5(n)=-1$ contribute destructively.
The result measures the asymmetry between split and inert sectors of the
Regge tower.

\smallskip\noindent\textbf{The complete circuit.}
The full determination of the odd zeta values forms a closed circuit
originating from $\golden^2 = \golden + 1$:
\[
\begin{tikzcd}[column sep=1.4em, row sep=2.2em]
  & \psi_p \ar[r] & \chi_5 \ar[r] \ar[d]
    & {\textstyle\prod f_p} \ar[r]
    & \zeta_K(s) \ar[r]
    & \zeta(2k{+}1) \ar[d] \\
  \golden^2{=}\golden{+}1 \ar[ur] \ar[dr]
    & & \mathrm{Li}_s(\omega^a) \ar[d]
    & & & a_{2k-1,0} \ar[d] \\
  & \omega{+}\omega^4{=}\golden^{-1} \ar[r]
    & L(s,\chi_5) \ar[uur, crossing over]
    & & & A_{\mathrm{Ven}} \ar[d] \\
  & & \text{Baker}(s{=}1) \ar[u, "\text{proved}"']
    & & & \mathcal{N}{=}4\ \text{SUSY}
\end{tikzcd}
\]
The upper branch (Frobenius $\to$ Euler product $\to$ Dedekind
factorisation) and the lower branch (Galois traces $\to$
polylogarithm $\to$ Baker) converge at the Veneziano amplitude.
The chain is unconditional through the structural determination
$\zeta(2k+1) = \zeta_{\mathbb{Q}(\sqrt{5})}(2k+1)/L(2k+1,\chi_5)$;
the remaining open questions---the algebraic independence of $\pi$ and $\log\phi$
and the transcendence of $L(2k+1,\chi_5)$ for $k\geq 1$---are
addressed in \cref{op:polylog-baker,op:periods}.
No cancellation mechanism internal to the
$\mathbb{Q}$-linear \cref{eq:polylog-rep}
distinguishes $s=1$ from $s\geq 3$.

The EFT of Elvang et al.\ operates in the additive world of
the Dirichlet series $\sum n^{-s}$: it sees the Regge tower as
an unstructured sum of states, with no access to the fundamental
theorem of arithmetic that converts the sum into a product and
reveals the $\chi_5$-grading.
The PCF arithmetic operates in the multiplicative fragment
(\cref{ssec:polylog-circuit}), where the
same values are structurally determined.
The Veneziano amplitude, singularised by $\mathcal{N}=4$ SUSY as
the unique UV completion, carries odd-zeta coefficients that are
fixed by the pentagonal arithmetic of $\mathbb{Q}(\sqrt{5})$:
SUSY selects the \emph{form} (Veneziano), the pentagon selects
the \emph{values} ($\zeta(2k+1)$), and no free parameter remains.

The isomorphism between $\golden$, one-dimensional real
arithmetic, and $\mathbb{R}$ (\cref{rmk:standard-id})
reveals a concrete relationship between the primes, the Euler
product, and $\zeta(s)$ on the one hand, and string theory on
the other.
A one-dimensional vector rotating inside a torus---a vibrating
string, a rotating vector in the
PCF framework---compactifies by accumulating winding numbers
(one per prime, via $\psi_p(\golden) = \golden^p$) and emerges
into two dimensions as it closes on the torus
$T^2_{\mathrm{PCF}}$.  In the PCF construction, it extends
further to three dimensions through $\golden$, which permits an
invariant relationship between 1D, 2D, and 3D: the golden coupling
$z = \golden y$ (\cref{prop:dim-chain}) extends
$\mathbb{C}$ to $E^3$, and the $S_3$ symmetry at modulus
$\mu_3 = 1/2$ (\cref{cor:S3-from-euler}) couples the
three eigenvalue directions to the complex plane, producing
invariants---$\ell = 1$, $G_N = 1/2$, $c = 3$---demonstrated
isomorphic to those of AdS/CFT~\cite{Maldacena,BrownHenneaux}
(\cref{lem:ads-radius}).
The algebraic content is concentrated in three properties of
$\golden^2 = \golden + 1$: the norm $N(\golden) = \golden\bar\golden = -1$
generates the S-duality that forces $\ell = 1$; the trace
$\operatorname{Tr}(\golden) = 1$ produces the self-duality of $\chi_5$
that closes the functional equation; and
$|\ker G| = 2$ contracts $\mu_2 = 1$ to $\mu_3 = 1/2$,
coupling the three eigenvalue directions via $S_3$.
A single algebraic relation manifests at every dimensional level.

These invariants, viewed from the perspective of Borger's~\cite{Borger}
Frobenius lifts ($\Lambda$-ring structure on $\Rpcf$), Manin's~\cite{Manin}
quantum tori programme, Lorscheid's~\cite{Lorscheid} torified varieties, and other
proposals for $\mathbb{F}_1$-geometry~\cite{ConnesDouglasSchwarz}---all descended from Weil's~\cite{Weil}
proof of the Riemann hypothesis for function fields over finite
fields, from which the programme to approach RH via absolute
geometry derives---are the foundations on which the Primitive
Structure~\cite{V11} builds its categorical framework.

What the present paper shows is how the Euler product, which
emerges from carrying the one-dimensional generator $\golden^\sigma$
to the sum over all primes, is isomorphic to the second dimension
expressed in $S_3$ symmetry, and how that dimension maintains an
invariant relationship with $E^3$.
The Euler product does not merely compute $\zeta(s)$: it
simultaneously assembles the second dimension
(convergence on~$\mathbb{C}$, upper branch of the diagram in
\cref{sec:objective}) and determines the spectral structure
that $S_3$ symmetry imposes on the third (lower branch).
The $S_3$ coupling between the eigenvalue directions $P$, $C$, $F$
at modulus $\mu_3 = 1/2$
(\cref{eq:tripartite-norm})
produces the same invariants from the arithmetic of $\golden$ alone.

The invariant relationship between dimensions is one of the reasons
the golden ratio has been essential in the art of diverse cultures
and epochs: the self-similarity $\golden^2 = \golden + 1$---a
rectangle subdividing into a square and a similar rectangle---is the
visual manifestation of the same dimensional invariance that,
iterated over the primes, produces the odd zeta values.
This ancient measurement principle allows the preservation of
structure across scales; what the Euler product computes is the
preservation of structure across primes.
Both are $\golden^2 = \golden + 1$.

This is why $\golden$ is the key to $\mathbb{F}_1$-descent without
breaking symmetry with the integers.  Descent to $\mathbb{F}_1$
strips addition; what survives is the multiplicative skeleton---powers,
norms, magnitudes.  Most generators would lose contact with the primes
at this stage: without addition, the Dirichlet series $\sum n^{-s}$
disappears, and with it the connection between the analytic world and
the arithmetic one.  But $\golden$ preserves the connection, because
the Frobenius lifts $\golden^p = F_p\golden + F_{p-1}$ encode the
splitting type of every prime $p$ through a \emph{purely
multiplicative} operation---the $p$-th power---whose residue modulo $p$
is $\chi_5(p)$ (\cref{prop:fib-split}).
The integers, and therefore the primes, remain visible at $\mathbb{F}_1$
because $\golden^\sigma\colon\mathbb{R}\xrightarrow{\sim}\mathbb{R}_{>0}$
is an isomorphism (\cref{rmk:standard-id}): the one-dimensional
vector $\golden$ \emph{is} $\mathbb{R}$ as a geometric object.
No other quadratic generator has this property
(\cref{prop:phi-unique}).

The critical step is the moment $\golden$ meets the circle---what
the Primitive Structure~\cite{V11} calls the \emph{Vitruvian modulus}.
The pentagon identity $\golden = 2\cos(\pi/5)$
(\cref{prop:pentagon}) embeds the self-similar generator
into the unit circle, and the torus $T^2_{\mathrm{PCF}}$ with
$\tau = i$ and period $T = 2\pi\log\golden$ completes the embedding.
The complex plane inherits from a dual genealogy---geometric
(Vitruvius, \emph{De Architectura}, c.~25~BC) and algebraic
(Cardano, \emph{Ars Magna}, 1545; Gauss, Argand~\cite{Argand})---the triple
$\{\text{origin, modulus, angle}\}$
that organises every complex number as $z = |z|\,e^{i\theta}$
(cf.\ the Primitive Structure~\cite{V11}, \S2).
In PCF, $\golden$ generates the origin through $\golden^2 = \golden + 1$,
$|\Omega| = 1/2$ is the modulus (radial magnitude),
and the $S_3$ eigenvalues $\omega^k/2$ provide the angular
directions at $120^\circ$ separation.
The one-dimensional relations that $\golden$ carries with respect to
itself---self-similarity, norm $-1$, trace $1$---are transported into
this two-dimensional structure without distortion: a
one-dimensional vector \emph{is} a two-dimensional circular manifold
($\mathbb{R} \xrightarrow{e^{i\theta}} S^1$), and the symmetries
it carries---periodicity, winding,
self-similarity---are preserved in higher dimensions, as
demonstrated in the Primitive Structure~\cite{V11} by the chain of
forgetful functors
$C^*$-$\mathbf{Alg}\to\mathbf{Hilb}\to\mathbf{Met}\to\mathbf{Set}$.
This is the geometric content of the Vitruvian principle: what
Brunelleschi formalised for perspective projection---encoding
three-dimensional structure in two dimensions while preserving the
ratios that define the object---$\golden$ accomplishes for the
descent from $\mathbb{C}$ to $\mathbb{F}_1$.
The multiplicative relations survive the loss of addition because
$\golden^2 = \golden + 1$ is its own projection formula.
The topological equivalence $\mathbb{R} \cong S^1$, equipped with
$\golden$ as the generator and iterated over the primes through
the Frobenius lifts, is the mechanism behind the odd zeta values.
The organising principle behind $\zeta(3)$, $\zeta(5)$, $\zeta(7)$,
\ldots\ is the dimensional invariance of $\golden^2 = \golden + 1$,
which allows for self-consistent dimensional emergence generating
the relationships between the PCF
framework and several string theory frameworks demonstrated in this
paper---through the Euler product.

Elvang et al.~\cite{Elvang2026} observe that the odd zeta
values appear to have no algebraic relationships with the even ones.
This is a consequence of the limitation identified above:
the EFT operates in the additive fragment, where the
Dedekind factorisation that connects even and odd values through
$\chi_5$ is invisible.
The relationship is architectural:
the compactification of the open string on a torus and the
compactification of $\golden^\sigma$ over the primes through
the Euler product are the same compactification.
The string vibrates over the primes through the Frobenius
lifts $\psi_p(\golden) = \golden^p = F_p\golden + F_{p-1}$;
each winding number $\golden^p$ is a return to the torus
weighted by Fibonacci.
The $\zeta(2k+1)$ that the Veneziano amplitude carries as
coefficients are the ratios
$\zeta_{\mathbb{Q}(\sqrt{5})}(2k+1)/L(2k+1,\chi_5)$ that this mechanism produces.

\begin{remark}[The diagonal obstruction and the free-parameter illusion]%
\label{rmk:lawvere-diagonal}
The apparent freedom of $\zeta(2k+1)$ in the EFT admits a precise
logical diagnosis.
Lawvere~\cite{Lawvere} proved that in any cartesian closed category
admitting a point-surjective morphism $f\colon A\times A\to Y$, every
endomorphism of~$Y$ has a fixed point; Yanofsky~\cite{Yanofsky} generalised this to
show that the paradoxes of Liar, Russell, G\"odel, and Turing are
instances of this single diagonal construction
$g(s) = \alpha(f(s,s))$.
The EFT possesses both additive and multiplicative structure
($\sum a_n\alpha^{\prime n}$ and the operator product): it lives in
a category with the diagonal morphism $f\colon A\times A\to Y$
required by Lawvere's theorem.  In such a setting, the values
$\zeta(2k+1)$ are analogous to G\"odel sentences: true (the
Veneziano amplitude uses them) but not derivable from within the
system.

The PCF construction escapes this obstruction by descending to the
multiplicative fragment $\Ztwenty$, where addition is destroyed
(Proposition~2 of~\cite{V11}, verified in Lean).  Without addition,
the diagonal $f(s,s)$ cannot be formed: the No-Diagonal theorem
(Theorem~3 of~\cite{V11}) establishes
$\|\Omega\| = 1/2 < 1$, blocking the point-surjectivity that
Lawvere's theorem requires.
This fragment is decidable in the sense of Tarski~\cite{tarski1951decision} (\cite{V11},
\S2.4): a purely multiplicative domain admits elimination of
quantifiers and is therefore complete.
Within this decidable fragment, $\zeta(2k+1)$ is determined
(\cref{thm:pentagon-odd}): the Frobenius lifts
$\psi_p(\golden) = \golden^p$, operating without addition,
classify every prime and assemble the Euler product that produces
$\zeta(2k+1) = \zeta_{\mathbb{Q}(\sqrt{5})}(2k+1)/L(2k+1,\chi_5)$.

The free parameter of the EFT is not a property of $\zeta(2k+1)$
but of the system from which it is observed: a system with the
diagonal structure sees indeterminacy; a system without it sees
determination.
\end{remark}

%%% ====================================================================
\subsection{Verification}\label{sec:verification}
%%% ====================================================================

All identities in the present work have been verified at 50-digit precision
using \texttt{mpmath} and \texttt{sympy} in the accompanying script
\texttt{verify\_odd\_zeta\_pcf\_unified.py}.  The 207 checks cover:
Frobenius lifts (\cref{sec:background}),
prime decomposition (\cref{sec:primes}),
Euler product convergence (\cref{sec:primes}),
factorisation (\cref{sec:factorisation}),
the fundamental formula (\cref{sec:fundamental}),
the even Bernoulli formula (\cref{sec:even}),
odd-value determination (\cref{sec:pentagon}),
and the categorical framework (\cref{sec:categorical-framework}):
Frobenius functoriality $\Psi(mn) = \Psi(m)\circ\Psi(n)$,
the four Dedekind factorisations of $\mathbf{QuadExt}_{20}$
(\cref{thm:dedekind-natural}),
partial Euler product convergence to the colimit
(\cref{prop:euler-colimit}), and
co-determination by~$G$ (\cref{thm:co-determination}),
terminality of $\mathrm{Spec}(\OmH)$
(\cref{def:SpecOmega}),
the $\Lambda$-ring congruence $F_p\equiv\chi_5(p)\pmod{p}$
(\cref{rmk:lambda-ring}),
closure of $\{1,\chi_4,\chi_5,\chi_4\chi_5\}$ under multiplication
(\cref{thm:dedekind-natural}),
and monoidal multiplicativity of partial Euler products
(\cref{rmk:monoidal}).
The classical bridge (\cref{rmk:classical-bridge}) is verified
through $G$-fiber consistency (same fiber implies same factor type),
the four Gelfand $L$-functions, and the spectral mapping theorem
applied to $\OmH^2$ and $\OmH^3$.
The Weil analogy (\cref{ssec:weil-analogy}) is verified through
the $G$-homomorphism (all $64$ pairs), surjectivity ($4$ fibers),
the upper bound (all $\mathrm{Re}(\mathrm{spec}(\OmH))\le\mu_3$),
the lower bound prerequisite (exponent~$2$), and the eigenvalue
moduli $|\omega^k/2|=1/2$.
The polylogarithmic circuit (\cref{ssec:polylog-circuit}) is
verified through the Galois traces $\omega+\omega^4 = 1/\golden$,
the Gauss sum $\tau(\chi_5) = \sqrt{5}$, the polylog \cref{eq:polylog-rep}
for $s = 3,5,7,9,11$, Baker's reduction at $s=1$
($|1-\omega^2|^2/|1-\omega|^2 = \golden^2$), the additive
destruction of $\Ztwenty$, and the Fourier projection
$\cos(2\pi n/5)-\cos(4\pi n/5) = (\sqrt{5}/2)\chi_5(n)$.
All checks pass with zero errors, including the Fibonacci splitting criterion
$\chi_5(p)\equiv F_p\pmod{p}$ (\cref{sec:frobenius-split}),
the Mersenne bridge $\golden^{\lambda_{\log}}=2$
(\cref{sec:rpcf-role}), the pentagonal formulas
(\cref{sec:pentagon}),
and the spectral structure: $\mu_3 = 1/2$, eigenvalues
$\omega^k/2$, $\ell = 1$ from S-duality, Brown--Henneaux
$c = 3$, and the diagonal obstruction
(\cref{ssec:unifying}, \cref{lem:ads-radius}).

The algebraic and spectral results of \cref{sec:bridges,ssec:unifying} have been formally verified in Lean~$4$: eigenvalue real parts
(\lean{eigenvalue_half}), spectral uniqueness
(\lean{spectral_uniqueness}), the Galois functor
(\lean{galois_functor_complete},~\cite{V11}), and the
pentagonal identity $\golden = 2\cos(\pi/5)$
(\lean{cos_pi_div_five_eq}, \lean{pentagon_identity}).

Table~\ref{tab:lean-audit} lists the correspondence between
the paper's results and the Lean identifiers in
\lean{odd_zeta_pcf_v1.lean}.

\begin{table}[ht]
\centering
\caption{Lean audit: correspondence between paper results and
formalised identifiers. All dependencies are fully verified.
The single axiom is $\golden^2=\golden+1$.}
\label{tab:lean-audit}
\footnotesize
\begin{tabular}{lll}
\toprule
Paper result & Lean identifier & Method \\
\midrule
\multicolumn{3}{l}{\textbf{Part 1: PCF isomorphism chain}} \\
$\golden^2=\golden+1$ & \lean{phi_sq} & axiom \\
$4(\golden/2)^2-2(\golden/2)-1=0$ & \lean{phi_half_quadratic} & \texttt{nlinarith} \\
$\cos(5\theta)$ expansion & \lean{cos_five_mul} & \texttt{nlinarith} \\
$\cos(\pi/5)>0$ & \lean{cos_pi_five_pos} & monotonicity \\
$4\cos^2(\pi/5)-2\cos(\pi/5)-1=0$ & \lean{cos_pi_five_quadratic} & unique root \\
$\golden = 2\cos(\pi/5)$ & \lean{pentagon_identity} & \texttt{linarith} \\
$|\Ztwenty|=8$ & \lean{Z20star_card} & \texttt{decide} \\
Mult.\ closed & \lean{mul_closed} & \texttt{decide} \\
Add.\ destroyed & \lean{add_destroyed} & \texttt{decide} \\
$G$ homomorphism & \lean{G_hom} & \texttt{decide} \\
$G$ surjective & \lean{G_surj} & \texttt{fin_cases} \\
$|\ker G|=2$ & \lean{G_kernel} & \texttt{decide} \\
Exponent~$2$ & \lean{exponent_two} & \texttt{decide} \\
$\mu_2/\mu_3 = 2 = |\ker G|$ & \lean{contraction_is_kernel} & \texttt{norm_num} \\
$(\sigma,\mu)=(3/2,1/2)$ unique & \lean{spectral_uniqueness} & \texttt{nlinarith} \\
Diagonal blocked & \lean{diagonal_blocked} & \texttt{nlinarith} \\
$\mathrm{Re}(\omega_{\mathrm{pcf}})=-1/2$ & \lean{$\omega$_properties} & \texttt{exp_mul_I} \\
$\mathrm{Re}>0\Rightarrow\mathrm{Re}=1/2$ & \lean{eigenvalue_half} & exhaustion \\
\midrule
\multicolumn{3}{l}{\textbf{Part 2: Isomorphism with Euler product}} \\
$\forall p>5$: $p\bmod 20\in\Ztwenty$ & \lean{primes_land_in_Z20star} & \texttt{omega+decide} \\
$\chi_5$ from $G$ & \lean{chi5_from_G} & \texttt{rfl} \\
$\chi_5$ classification & \lean{chi5_values} & \texttt{decide} \\
$\chi_5(p)\equiv F_p\pmod{p}$ & \lean{fib_split_criterion} & \texttt{native_decide} \\
$f_p$ type from $G$ & \lean{euler_type_from_G} & \texttt{rfl} \\
\midrule
\multicolumn{3}{l}{\textbf{Part 3: Transport}} \\
$\zeta$ spectral $=$ PCF & \lean{spectral_identification} & \texttt{rfl} \\
Transport principle & \lean{transport_spectral} & rewrite \\
$\zeta$ positive $\Rightarrow 1/2$ & \lean{zeta_positive_spectral_half} & transport \\
\midrule
\multicolumn{3}{l}{\textbf{Part 4: $\mathbb{F}_1$-descent, Weil, physics}} \\
$z^2 = zy + y^2$ & \lean{self_measurement} & \texttt{nlinarith} \\
$\psi_p\circ\psi_q = \psi_{pq}$ & \lean{frobenius_composition} & \texttt{pow_mul} \\
$\mathbb{F}_1$-descent data & \lean{F1_descent_data} & conjunction \\
Weil analogy & \lean{weil_analogy} & conjunction \\
$\mu_3^2 = 1/4$ & \lean{mu3_squared} & \texttt{norm_num} \\
$1-\mu_3^2 = 3/4$ & \lean{holographic_ratio} & \texttt{norm_num} \\
$c = 3\ell/(2G_N) = 3$ & \lean{brown_henneaux_c} & \texttt{norm_num} \\
$1/(4G_N) = \mu_3$ & \lean{BH_factor} & \texttt{norm_num} \\
$-1/i = i$ (S-duality) & \lean{S_duality_fixed} & \texttt{Complex.I_mul_I} \\
$\ell=1/\ell \Rightarrow \ell=1$ & \lean{ads_radius_one} & \texttt{nlinarith} \\
$\bar\golden = 1-\golden$ & \lean{phi_bar_eq} & \texttt{field_simp} \\
$\mathrm{Tr}(\golden) = 1$ & \lean{trace_one} & \texttt{ring} \\
$N(\golden) = -1$ & \lean{norm_neg_one} & \texttt{field_simp} \\
$\mathrm{disc}(\mathbb{Q}(\sqrt{5})) = 5$ & \lean{discriminant_five} & \texttt{nlinarith} \\
\bottomrule
\end{tabular}
\end{table}

The categorical structures of \cref{sec:categorical-framework}---the
Frobenius functor (functoriality verified via
$\psi_p\circ\psi_q=\psi_{pq}$), the co-determination theorem
(reducing to \lean{G_hom}, \lean{G_surj},
\lean{G_kernel} in \lean{odd_zeta_pcf_v1.lean};
equivalently \lean{G_hom}, \lean{G_surj},
\lean{G_kernel} in~\cite{V11}), and the spectral uniqueness
(\lean{spectral_uniqueness})---are
verified through the Lean~$4$ formalisation of the Primitive
Structure~\cite{V11}.

%%% ====================================================================
\subsection{Open problems}\label{sec:open}
%%% ====================================================================

\begin{openproblem}[Odd Zeta Irrationality]\label{op:irrationality}
Does the pentagonal Hurwitz representation (\cref{thm:hurwitz}) permit
a modular proof of the irrationality of $\zeta(2k+1)$ for $k \ge 2$?
The case $k=1$ (Ap\'ery) is known, but the higher cases lack a
general proof.
\end{openproblem}

\begin{openproblem}[PCF Modularity]\label{op:modularity}
Is the PCF torus $T^2_{\mathrm{PCF}}$ at $\tau=i$ a component of a
larger modular object that explains the choice of $\chi_5$ over
other quadratic characters?
\end{openproblem}

\begin{openproblem}[Physical Formalization]\label{op:physical}
Can the holographic identity $S_{\mathrm{BH}} \sim L(1,\chi_5)$ be
extended to Higher-Spin theories or String Field Theory using the
full $S_3$ symmetry?
\end{openproblem}

\begin{openproblem}[F1-descent]\label{op:f1}
Formalize the full decomposition $\zeta = \zeta_{\mathbb{Q}(\sqrt{5})}/L$
in terms of Borger's $\Lambda$-rings within Lean 4.
\end{openproblem}

\begin{openproblem}[Galois Functor Surjectivity]\label{op:surjectivity}
Does the character group of $\mathbb{Z}_2 \times \mathbb{Z}_2$
satisfy a categorical limit theorem that forces the distribution of
primes in $\chi_d$ fibers?
\end{openproblem}

%%% ====================================================================
\subsection{Conclusion}\label{sec:conclusion}
%%% ====================================================================

The chain from the pentagon to the odd zeta values
(\cref{prop:pentagon-dictates,rmk:pentagonal-dictate}):
\[
\begin{split}
  &\golden=2\cos\tfrac{\pi}{5}
  \;\to\;
  \chi_5(a)=\tfrac{\log|2\cos(a\pi/5)|}{\log\golden}\\[3pt]
  &\quad\to\;
  L(s,\chi_5)=\tfrac{\textstyle\sum_a\chi_5(a)\,\zeta(s,a/5)}{5^s}
  \;\to\;
  \zeta(2k+1)=\tfrac{\zeta_{\mathbb{Q}(\sqrt{5})}(2k+1)}{L(2k+1,\chi_5)},
\end{split}
\]
shows that the odd zeta values $\zeta(2k+1)$ are structurally
determined by $\golden$ and $\pi$ via the pentagonal arithmetic
of $\mathbb{Q}(\sqrt{5})$.  Every step in this chain derives from
the single relation $\golden^2 = \golden + 1$
(\cref{thm:unifying}).

The one-dimensional generator $\golden^\sigma$ reaches the
two-dimensional torus $T^2_{\mathrm{PCF}}$ by two routes:
geometrically, the tower closes with period $T = 2\pi R$
($R = \log\golden$); arithmetically, the Frobenius lifts
$\golden^p = F_p\golden + F_{p-1}$, iterated over all primes,
assemble the Euler product $\zeta_{\mathbb{Q}(\sqrt{5})}(s)$ on~$\mathbb{C}$.
The dimensional structure emerges deterministically
(\cref{cor:S3-from-euler}):
\[
  \underbrace{\golden^2 = \golden+1}_{\textup{1D generator}}
  \;\to\;
  \underbrace{\prod_p f_p(s)}_{\textup{Euler product on }\mathbb{C}}
  \;\to\;
  \underbrace{|\ker G|=2}_{\textup{contraction}}
  \;\to\;
  \underbrace{\mu_3 = \tfrac{1}{2}}_{\textup{norm}}
  \;\to\;
  \underbrace{S_3}_{\textup{symmetry}}
  \;\to\;
  \underbrace{E^3,\; c=3.}_{\textup{geometry}}
\]
No dimension is postulated; each is derived from $\golden^2 = \golden + 1$.

The regulator $R = \log\golden$ is the common invariant:
\begin{itemize}
\item the \textbf{arithmetic size} of the field: $L(1,\chi_5)=2R/\sqrt{5}$;
\item the \textbf{worldsheet period}: $T_{\mathrm{PCF}}=2\pi R$ (\cref{def:torus});
\item the \textbf{entropy anchor}: $S_{\mathrm{BH}}/k_B=\lambda_{\log}\cdot R$ (\cref{prop:worldline-entropy}).
\end{itemize}
The categorical framework of \cref{sec:categorical-framework}
makes the structure of this determination precise: the Frobenius
system is a functor from the multiplicative monoid of integers to
the endomorphisms of $\Rpcf$; the Euler product is the colimit of
the partial products determined by $G$; the Dedekind factorisation
is a natural isomorphism across the four quadratic extensions
classified by $G$; and the two branches---arithmetic (values of
$\zeta$) and spectral (eigenvalues of $\OmH$)---form a span with
common source $G$.  This span provides the mathematical content of
the \texttt{bridge\_axiom} formalised in the Primitive
Structure~\cite{V11}: the single map $G$ that produces the odd zeta
values also produces the spectral constraint $\mu_3=1/2$.

The present paper provides both the arithmetic reading of
$T^2_{\mathrm{PCF}}$ (class-number formula, Euler product,
Frobenius structure) and its physical connections
(partition function, Bekenstein--Hawking entropy,
Brown--Henneaux central charge $c = 3$).
The values $\zeta(2k+1)$ that appear free in the $\mathcal{N}=4$
EFT of Elvang et al.~\cite{Elvang2026} are dictated by the pentagon:
the EFT identifies the unique UV completion (Veneziano); PCF identifies
why its odd-zeta parameters take the values they do.

